\chapter{Approximation and Density}
\label{ch:iii22}

Approximation and density justify replacing rough functions by smooth ones. Mollifiers, cutoffs, and extension arguments let proofs and numerical constructions move between Sobolev or $L^p$ functions and smooth compactly supported models.

\section*{Learning goals}
After this chapter, the reader should be able to:
\begin{itemize}
\item construct smooth approximations by mollification and quantify the mode of convergence;
\item prove density of smooth or compactly supported smooth functions in the function spaces covered by the chapter;
\item control boundary effects when approximating functions on domains rather than on all of \(\mathbb R^n\);
\item use cutoff functions and convolution in the correct order to preserve desired support properties;
\item distinguish approximation in norm from pointwise approximation;
\item identify spaces or boundary conditions for which naive mollification does not directly give the required dense subclass;
\end{itemize}
\section*{Conceptual roadmap}
\[
\boxed{\text{structure}\;\longrightarrow\;\text{estimate}\;\longrightarrow\;\text{limit or representation}\;\longrightarrow\;\text{application}.}
\]

\paragraph{Mathematical setting.}
Convergence statements in $L^p$ and $W^{k,p}$ below use $1\le p<\infty$. On an arbitrary open $\Omega$, convolution identities and Sobolev convergence hold on each $\Omega'\Subset\Omega$ for sufficiently small mollifier support. On $\mathbb R^n$ they hold globally. Standard density/extension theorems, with full proofs deferred: $C^\infty(\Omega)\cap W^{k,p}(\Omega)$ is dense in $W^{k,p}(\Omega)$ on any open set; compactly supported tests are dense in $W_0^{1,p}$ by definition, not in all $W^{1,p}$. A bounded Lipschitz domain admits a bounded extension $W^{1,p}(\Omega)\to W^{1,p}(\mathbb R^n)$. For the frequency example, $H^s(\mathbb R^n)$ uses the weight $(1+4\pi^2|\xi|^2)^s$ and $\chi_R(\xi)=\chi(\xi/R)$ with one fixed smooth compactly supported $\chi=1$ near zero.

\section{Mollifiers}
A standard mollifier is smooth, nonnegative, compactly supported, and normalized to integral one.


% BEGIN VOL03-EXPANSION III22-example-01
\begin{example}[Mollifying a discontinuous indicator]\label{ex:iii22-expand-01}
Let \(f=1_{[0,1]}\) on \(\mathbb R\) and let \(\rho_\varepsilon\) be a standard
mollifier. Then
\[
f_\varepsilon=f*\rho_\varepsilon
\]
is smooth, equals \(1\) on \([\varepsilon,1-\varepsilon]\), equals zero outside \([-\varepsilon,1+\varepsilon]\), and and transitions smoothly near \(0\) and \(1\). For every finite
\(p\),
\[
\|f_\varepsilon-f\|_p\to0.
\]
The discrepancy is confined to shrinking boundary layers.
\end{example}
% END VOL03-EXPANSION III22-example-01
\section{Approximate identities}
Scaled mollifiers concentrate at the origin and approximate functions in $L^p$.

\section{Smooth approximation in $L^p$}
Convolution plus cutoff proves density of $C_c^\infty$ in $L^p$ for $1\le p<\infty$.

\section{Sobolev approximation}
Mollification approximates weak derivatives simultaneously on interior subdomains.

\section{Boundary issues}
Near the boundary, extension operators or geometry-aware cutoffs are needed to mollify without leaving the domain.


% BEGIN VOL03-EXPANSION III22-example-02
\begin{example}[Why boundary-aware extension is needed]\label{ex:iii22-expand-02}
If \(f\) is defined only on a domain \(\Omega\), direct convolution samples
values \(f(x-y)\) that may lie outside \(\Omega\) when \(x\) is near the
boundary. One therefore extends \(f\) to a larger region or first localizes
with cutoffs. This is not a technical nuisance: without an extension rule,
the convolution is not even defined from the original data near the boundary.
\end{example}
% END VOL03-EXPANSION III22-example-02
\section{Density of test functions}
$C_c^\infty(\Omega)$ is dense in $W_0^{1,p}(\Omega)$ by definition and under standard characterizations.

\section{Frequency truncation}
Fourier cutoffs give an alternative smooth approximation method in whole-space Sobolev spaces.


% BEGIN VOL03-EXPANSION III22-example-03
\begin{example}[Frequency cutoff smooths whole-space Sobolev data]\label{ex:iii22-expand-03}
For \(u\in H^s(\mathbb R^n)\), choose a smooth frequency cutoff
\(\chi_R(\xi)\) supported in \(|\xi|\lesssim R\) and equal to one on a smaller
ball. Define
\[
\widehat u_R=\chi_R\widehat u.
\]
Then \(u_R\) is smooth and
\[
\|u_R-u\|_{H^s}\to0
\]
by dominated convergence with the Sobolev weight. Frequency truncation is the
spectral analogue of mollification.
\end{example}
% END VOL03-EXPANSION III22-example-03
\section{Approximation consequences}
Density transfers identities and inequalities proved for smooth functions to rough functions by norm limits.

\section{Core structural results}

\begin{theorem}[$L^p$ mollifier convergence]\label{thm:iii22-01}
If $f\in L^p(\mathbb R^n)$, $1\le p<\infty$, then $f*\rho_\varepsilon\to f$ in $L^p$.
\begin{proof}
Write the convolution difference as an average of translations $\tau_yf-f$ and use continuity of translations in $L^p$ together with concentration of the mollifier near zero.
\end{proof}
\end{theorem}

\begin{theorem}[Derivative commutation]\label{thm:iii22-02}
If $f\in W^{k,p}(\Omega)$, then $D^\alpha(f*\rho_\varepsilon)=(D^\alpha f)*\rho_\varepsilon$ for $|\alpha|\le k$ at points with distance greater than $\varepsilon$ from the boundary. On $\mathbb R^n$ the identity is global.
\begin{proof}
Use the weak derivative identity to move derivatives between the function and mollifier, then apply convolution associativity on compactly supported kernels.
\end{proof}
\end{theorem}

\begin{theorem}[Density of $C_c^\infty$ in $L^p$]\label{thm:iii22-03}
For $1\le p<\infty$, smooth compactly supported functions are dense in $L^p(\mathbb R^n)$.
\begin{proof}
First truncate $f$ spatially and in magnitude to obtain bounded finite-support approximants, then mollify them; each step has arbitrarily small $L^p$ error.
\end{proof}
\end{theorem}

\section{Worked examples}

\begin{example}[Mollifying a rough function]\label{ex:iii22-worked-01}
Let \(u=\mathbf1_{(0,1)}\) and let \(\rho_\varepsilon\) be a standard
mollifier. Then
\[
u_\varepsilon=u*\rho_\varepsilon
\]
is smooth, agrees with values near \(1\) away from the jump layers, and
converges to \(u\) in \(L^p_{\mathrm{loc}}\) for finite \(p\). The jump is
replaced by a transition layer of width \(O(\varepsilon)\).
\end{example}

\begin{example}[Derivative commutes with mollification]\label{ex:iii22-worked-02}
If \(u\) has weak derivative \(D^\alpha u\), then away from boundary
issues,
\[
D^\alpha(u*\rho_\varepsilon)
=(D^\alpha u)*\rho_\varepsilon
=u*(D^\alpha\rho_\varepsilon).
\]
Thus mollification approximates both a Sobolev function and its weak
derivatives.
\end{example}

\begin{example}[Boundary extension is a real issue]\label{ex:iii22-worked-03}
On a bounded domain, the formula \(u*\rho_\varepsilon\) requires values of
\(u\) outside \(\Omega\). Zero extension is harmless for \(H_0^1\) under the
standard hypotheses, but an arbitrary \(H^1\) function may develop an
artificial boundary jump. Density arguments must therefore specify an
extension or work on shrunken interior domains.
\end{example}

\begin{example}[Smooth functions are dense in \(L^p\)]\label{ex:iii22-worked-04}
For \(1\le p<\infty\), first truncate/localize an \(L^p\) function, then
convolve with a mollifier. Translation continuity in \(L^p\) yields
\[
\|u*\rho_\varepsilon-u\|_p\to0.
\]
The two steps handle tails and local smoothing separately.
\end{example}

\section{Solved dossiers}
Each dossier combines several ideas from the chapter in a longer argument. The aim is to make hypotheses, calculations, and proof strategy explicit.

\begin{problem}[Mollifying a rough function]\label{prob:iii22-dossier-01}
Describe what mollification does to the indicator of an interval.
\end{problem}
\begin{solution}
Convolution with a smooth approximate identity produces a smooth transition
near the endpoints and converges back to the indicator in \(L^p\) for finite
\(p\).
\end{solution}

\begin{problem}[Derivative commutes with mollification]\label{prob:iii22-dossier-02}
Show weak differentiation commutes with convolution by a mollifier.
\end{problem}
\begin{solution}
Move the derivative in the distributional pairing, or differentiate the
smooth kernel; both yield the same convolution identity.
\end{solution}

\begin{problem}[Boundary extension is a real issue]\label{prob:iii22-dossier-03}
Why can naive zero extension spoil Sobolev regularity at a boundary?
\end{problem}
\begin{solution}
A nonzero trace becomes a jump across the boundary after zero extension,
creating a singular derivative. One needs an extension theorem or zero-trace
hypothesis.
\end{solution}

\begin{problem}[Smooth functions are dense in \(L^p\)]\label{prob:iii22-dossier-04}
Outline a proof that \(C_c^\infty(\mathbb R^n)\) is dense in \(L^p(\mathbb R^n)\) for finite \(p\).
\end{problem}
\begin{solution}
First cut off the tails, then mollify the compactly supported approximation;
both errors can be made arbitrarily small in \(L^p\).
\end{solution}

\begin{problem}[Approximate identity in \(L^p\)]\label{prob:iii22-dossier-05}
State the basic \(L^p\) convergence of mollifiers.
\end{problem}
\begin{solution}
For \(1\le p<\infty\), \(u*\rho_\varepsilon\to u\) in \(L^p\) as
\(\varepsilon\downarrow0\).
\end{solution}

\begin{problem}[Cutoff localization]\label{prob:iii22-dossier-06}
If \(\chi_R\) is a smooth cutoff tending to \(1\), why does \(\chi_Ru\to u\) in \(L^p\)?
\end{problem}
\begin{solution}
The error is supported in the tail where \(u\) has arbitrarily small
\(L^p\) mass; dominated or absolute continuity of the integral completes the
argument.
\end{solution}

\begin{problem}[Sobolev density]\label{prob:iii22-dossier-07}
What must be controlled to prove smooth density in \(W^{1,p}\)?
\end{problem}
\begin{solution}
Both the function and each weak derivative must converge in \(L^p\).
\end{solution}

\begin{problem}[Compact support after mollification]\label{prob:iii22-dossier-08}
Why can mollification enlarge support?
\end{problem}
\begin{solution}
Convolution support lies in the Minkowski sum of the original support and
the mollifier support, enlarging it by roughly \(\varepsilon\).
\end{solution}

\begin{problem}[$L^p$ mollifier convergence proof dossier]\label{prob:iii22-dossier-09}
Prove the following structural statement and identify the step where its main hypothesis is used: If $f\in L^p(\mathbb R^n)$, $1\le p<\infty$, then $f*\rho_\varepsilon\to f$ in $L^p$.
\end{problem}
\begin{solution}
Write the convolution difference as an average of translations $\tau_yf-f$ and use continuity of translations in $L^p$ together with concentration of the mollifier near zero. This proof also explains why the stated hypothesis is structural rather than cosmetic.
\end{solution}

\begin{problem}[Derivative commutation proof dossier]\label{prob:iii22-dossier-10}
Prove the following structural statement and identify the step where its main hypothesis is used: If $f\in W^{k,p}(\Omega)$, then $D^\alpha(f*\rho_\varepsilon)=(D^\alpha f)*\rho_\varepsilon$ for $|\alpha|\le k$ at points with distance greater than $\varepsilon$ from the boundary. On $\mathbb R^n$ the identity is global.
\end{problem}
\begin{solution}
Use the weak derivative identity to move derivatives between the function and mollifier, then apply convolution associativity on compactly supported kernels. This proof also explains why the stated hypothesis is structural rather than cosmetic.
\end{solution}

\begin{problem}[Density of $C_c^\infty$ in $L^p$ proof dossier]\label{prob:iii22-dossier-11}
Prove the following structural statement and identify the step where its main hypothesis is used: For $1\le p<\infty$, smooth compactly supported functions are dense in $L^p(\mathbb R^n)$.
\end{problem}
\begin{solution}
First truncate $f$ spatially and in magnitude to obtain bounded finite-support approximants, then mollify them; each step has arbitrarily small $L^p$ error. This proof also explains why the stated hypothesis is structural rather than cosmetic.
\end{solution}

\begin{problem}[Chapter synthesis]\label{prob:iii22-dossier-12}
Connect the main definitions and structural theorems of this chapter into one proof strategy. Explain what object is controlled, what estimate or closure principle is used, and what conclusion becomes available.
\end{problem}
\begin{solution}
Approximation and density justify replacing rough functions by smooth ones. Mollifiers, cutoffs, and extension arguments let proofs and numerical constructions move between Sobolev or $L^p$ functions and smooth compactly supported models. The recurring strategy is to encode the object in the chapter's natural measurable, normed, Fourier, or distributional structure, establish the controlling estimate, and then pass to the desired limit or representation.
\end{solution}
\section{Exercises with complete solutions}
The shorter exercise layer remains separate from the solved dossiers.

\begin{exercise}\label{exr:iii22-01}
State and apply the central principle behind Mollifiers. Give one immediate consequence in the setting of this chapter.
\end{exercise}
\begin{hint}
Extend or localize the function first if naive convolution would sample outside the domain.
\end{hint}
\begin{solution}
A standard mollifier is smooth, nonnegative, compactly supported, and normalized to integral one. As an immediate consequence, the same principle can be inserted into later convergence, approximation, transform, or weak-form arguments whenever its hypotheses are verified.
\end{solution}

\begin{exercise}\label{exr:iii22-02}
State and apply the central principle behind Approximate identities. Give one immediate consequence in the setting of this chapter.
\end{exercise}
\begin{hint}
Choose a mollifier \(\rho_\varepsilon\) and use \(f*\rho_\varepsilon\) as the smooth approximation.
\end{hint}
\begin{solution}
Scaled mollifiers concentrate at the origin and approximate functions in $L^p$. As an immediate consequence, the same principle can be inserted into later convergence, approximation, transform, or weak-form arguments whenever its hypotheses are verified.
\end{solution}

\begin{exercise}\label{exr:iii22-03}
State and apply the central principle behind Smooth approximation in $L^p$. Give one immediate consequence in the setting of this chapter.
\end{exercise}
\begin{hint}
Norm convergence follows by translation continuity plus averaging; write the convolution error as an average of translates.
\end{hint}
\begin{solution}
Convolution plus cutoff proves density of $C_c^\infty$ in $L^p$ for $1\le p<\infty$. As an immediate consequence, the same principle can be inserted into later convergence, approximation, transform, or weak-form arguments whenever its hypotheses are verified.
\end{solution}

\begin{exercise}\label{exr:iii22-04}
State and apply the central principle behind Sobolev approximation. Give one immediate consequence in the setting of this chapter.
\end{exercise}
\begin{hint}
Use a cutoff after mollification when compact support inside the domain is required.
\end{hint}
\begin{solution}
Mollification approximates weak derivatives simultaneously on interior subdomains. As an immediate consequence, the same principle can be inserted into later convergence, approximation, transform, or weak-form arguments whenever its hypotheses are verified.
\end{solution}

\begin{exercise}\label{exr:iii22-05}
State and apply the central principle behind Boundary issues. Give one immediate consequence in the setting of this chapter.
\end{exercise}
\begin{hint}
Track weak derivatives through convolution: \(D^\alpha(f*\rho_\varepsilon)=(D^\alpha f)*\rho_\varepsilon\) in the weak sense.
\end{hint}
\begin{solution}
Near the boundary, extension operators or geometry-aware cutoffs are needed to mollify without leaving the domain. As an immediate consequence, the same principle can be inserted into later convergence, approximation, transform, or weak-form arguments whenever its hypotheses are verified.
\end{solution}

\begin{exercise}\label{exr:iii22-06}
State and apply the central principle behind Density of test functions. Give one immediate consequence in the setting of this chapter.
\end{exercise}
\begin{hint}
Boundary conditions may not survive ordinary mollification; modify the construction rather than assuming they do.
\end{hint}
\begin{solution}
$C_c^\infty(\Omega)$ is dense in $W_0^{1,p}(\Omega)$ by definition and under standard characterizations. As an immediate consequence, the same principle can be inserted into later convergence, approximation, transform, or weak-form arguments whenever its hypotheses are verified.
\end{solution}

\begin{exercise}\label{exr:iii22-07}
State and apply the central principle behind Frequency truncation. Give one immediate consequence in the setting of this chapter.
\end{exercise}
\begin{hint}
Density in norm is stronger than pointwise approximation; estimate the actual target-space norm.
\end{hint}
\begin{solution}
Fourier cutoffs give an alternative smooth approximation method in whole-space Sobolev spaces. As an immediate consequence, the same principle can be inserted into later convergence, approximation, transform, or weak-form arguments whenever its hypotheses are verified.
\end{solution}

\begin{exercise}\label{exr:iii22-08}
State and apply the central principle behind Approximation consequences. Give one immediate consequence in the setting of this chapter.
\end{exercise}
\begin{hint}
Choose the approximation scale after fixing the error tolerance so each localization and smoothing error is controlled separately.
\end{hint}
\begin{solution}
Density transfers identities and inequalities proved for smooth functions to rough functions by norm limits. As an immediate consequence, the same principle can be inserted into later convergence, approximation, transform, or weak-form arguments whenever its hypotheses are verified.
\end{solution}


% BEGIN VOL03-EXPANSION III22-exercises-01
\section{Graded supplementary exercises}
These exercises extend the preserved foundational set and are graded by mathematical role.

\subsection*{Standard computations}

\begin{exercise}[Mollifier mass]\label{exr:iii22-expand-01}
What is the integral of rho\_epsilon?
\end{exercise}
\begin{hint}
Scale variables.
\end{hint}
\begin{solution}
1.
\end{solution}

\begin{exercise}[Mollifier support]\label{exr:iii22-expand-02}
If rho is supported in the unit ball, where is rho\_epsilon supported?
\end{exercise}
\begin{hint}
Scale by epsilon.
\end{hint}
\begin{solution}
In the ball of radius epsilon.
\end{solution}

\begin{exercise}[Smoothness]\label{exr:iii22-expand-03}
Why is f*rho\_epsilon smooth for locally integrable f?
\end{exercise}
\begin{hint}
Derivatives fall on rho\_epsilon.
\end{hint}
\begin{solution}
All derivatives are given by convolution with smooth compactly supported kernel derivatives.
\end{solution}

\begin{exercise}[Lp convergence]\label{exr:iii22-expand-04}
For finite p, what happens to f*rho\_epsilon for f in Lp?
\end{exercise}
\begin{hint}
Use approximate identity theorem.
\end{hint}
\begin{solution}
It converges to f in Lp.
\end{solution}

\begin{exercise}[Weak derivative]\label{exr:iii22-expand-05}
If f is W1p, how does derivative of its mollification relate to Df?
\end{exercise}
\begin{hint}
Use commutation.
\end{hint}
\begin{solution}
D(f*rho\_epsilon)=(Df)*rho\_epsilon.
\end{solution}

\subsection*{Proofs}

\begin{exercise}[Approximate identity proof]\label{exr:iii22-expand-06}
Outline Lp mollifier convergence.
\end{exercise}
\begin{hint}
Write the difference as an average of translations.
\end{hint}
\begin{solution}
Translation continuity and kernel concentration give the result.
\end{solution}

\begin{exercise}[Derivative commutation]\label{exr:iii22-expand-07}
Prove weak derivatives commute with mollification.
\end{exercise}
\begin{hint}
Transfer derivatives using the weak identity.
\end{hint}
\begin{solution}
The smooth derivative equals convolution with the weak derivative.
\end{solution}

\begin{exercise}[Density]\label{exr:iii22-expand-08}
Explain the cutoff-plus-mollifier proof of density of Cc-infinity in Lp.
\end{exercise}
\begin{hint}
First truncate spatially, then mollify.
\end{hint}
\begin{solution}
Each step contributes arbitrarily small norm error.
\end{solution}

\begin{exercise}[Boundary need]\label{exr:iii22-expand-09}
Why does naive zero extension sometimes fail to preserve W1p regularity?
\end{exercise}
\begin{hint}
A nonzero trace creates a jump across the boundary.
\end{hint}
\begin{solution}
That jump contributes a singular derivative; extension operators or zero-trace hypotheses are needed.
\end{solution}

\subsection*{Counterexamples and hypothesis testing}

\begin{exercise}[Linfinity density]\label{exr:iii22-expand-10}
Why is Cc-infinity not dense in Linfinity on the whole line?
\end{exercise}
\begin{hint}
Uniform norm cannot approximate a nonvanishing constant by compact support.
\end{hint}
\begin{solution}
Outside every compact support the error from constant 1 is 1.
\end{solution}

\begin{exercise}[Pointwise not norm]\label{exr:iii22-expand-11}
Why does pointwise convergence of mollifications not alone prove Lp density?
\end{exercise}
\begin{hint}
Norm convergence requires integrated error control.
\end{hint}
\begin{solution}
Approximate-identity estimates supply that stronger control.
\end{solution}

\begin{exercise}[Boundary leakage]\label{exr:iii22-expand-12}
Why can direct mollification violate domain support or boundary conditions?
\end{exercise}
\begin{hint}
Convolution spreads support by the kernel radius.
\end{hint}
\begin{solution}
Values can cross the boundary unless cutoffs/extensions are used.
\end{solution}

\subsection*{Applications and investigations}

\begin{exercise}[Finite elements]\label{exr:iii22-expand-13}
Why is density important in variational numerics?
\end{exercise}
\begin{hint}
Smooth/finite-dimensional test functions approximate energy-space functions.
\end{hint}
\begin{solution}
Identities proved on simple functions extend to weak solutions by norm limits.
\end{solution}

\begin{exercise}[Spectral regularization]\label{exr:iii22-expand-14}
Why use frequency cutoff rather than physical mollification in whole-space PDE?
\end{exercise}
\begin{hint}
It diagonalizes Sobolev norms and operator multipliers.
\end{hint}
\begin{solution}
Regularity and convergence can be read directly from weighted spectral tails.
\end{solution}

\subsection*{Challenge problems}

\begin{exercise}[Two-scale approximation]\label{exr:iii22-expand-15}
Describe how to approximate an Lp function by a smooth compactly supported one within epsilon.
\end{exercise}
\begin{hint}
Choose a compact cutoff making the tail small, then choose mollifier scale making local smoothing error small.
\end{hint}
\begin{solution}
Triangle inequality combines the two errors below epsilon.
\end{solution}

\begin{exercise}[Sobolev density transfer]\label{exr:iii22-expand-16}
Why does derivative commutation allow mollifier convergence in Wkp rather than just Lp?
\end{exercise}
\begin{hint}
Apply the Lp approximate-identity theorem to every weak derivative.
\end{hint}
\begin{solution}
Each derivative component converges, so the summed Sobolev norm converges.
\end{solution}
% END VOL03-EXPANSION III22-exercises-01
\section*{Chapter summary}
The chapter now supplies a canonical theorem-and-dossier layer for subsequent measure, Fourier, distributional, Sobolev, and PDE arguments.
